% ============================================================
% Overseas Pipeline -- Cover Letter Template
%
% 使用方法（由 Claude Code 在 Step 3 生成时替换占位符）：
%
% 占位符说明：
%   {{POSITION_TITLE}}     - 职位名称，如 "Lecturer (Level B)"
%   {{POSITION_REF}}       - 职位 Reference 号（如有），如 "Ref: 688354"
%   {{DEPARTMENT_NAME}}    - 院系名称，如 "Department of Data Science and AI"
%   {{UNIVERSITY_NAME}}    - 学校名称，如 "Monash University"
%   {{SEARCHING_FOCUS}}    - 该职位的研究重点描述
%   {{FACULTY_OVERLAP}}    - 与 faculty 的具体 overlap 段落（1-2 句）
%   {{CLOSING_PARA}}       - 结尾段落（感谢 + 期望讨论）
%   {{JOB_URL}}            - 职位链接
%   {{GRANT_MENTION}}      - 资助计划提及（如澳洲的 ARC DECRA）
%
% 编译：xelatex main.tex
% ============================================================

\documentclass[11pt]{OUCletter}
\usepackage{fontspec}
\usepackage{tikz}
\usepackage{xcolor}
\usepackage{hyperref}

\hypersetup{
  colorlinks=true,
  linkcolor=black,
  citecolor=blue,
  urlcolor=blue,
  filecolor=magenta,
  pdfborder={0 0 0},
  breaklinks=true,
  bookmarksopen=true,
  bookmarksnumbered=true
}

\definecolor{scured}{RGB}{1, 45, 114}
\usepackage{fancyhdr}
\usepackage{lastpage}
\usepackage{eso-pic}

\pagestyle{fancy}
\fancyhf{}
\renewcommand{\headrulewidth}{0pt}
\renewcommand{\footrulewidth}{0pt}
\rhead{Page \thepage \hspace{1pt} of \pageref{LastPage}}

% ============================================================
% 自定义变量 -- Claude Code 在生成时替换这些值
% ============================================================
\newcommand{\searchingfocus}{\textbf{{{SEARCHING_FOCUS}}}}
\newcommand{\departmentname}{{{DEPARTMENT_NAME}}}
\newcommand{\universityname}{{{UNIVERSITY_NAME}}}

% ============================================================
% 申请人信息（Sophia 固定信息，通常不需要修改）
% ============================================================
\def\name{\textbf{Jie Gao} \\[10pt]
    Malone Postdoctoral Fellow \\
    Department of Computer Science \\
    Malone Center for Engineering in Healthcare \\
    Center for Language and Speech Processing \\
    Johns Hopkins University \\
    jgao77@jh.edu \\
    \href{https://gaojie058.github.io/}{https://gaojie058.github.io/} \\
}

\def\Where{\hspace{-1.2mm}\textbf{\color{scured}%
        {Department of Computer Science}, \\
        {Johns Hopkins University}
}}

\def\Address{Malone Hall, 3400 N Charles St}
\def\CityZip{Baltimore, MD 21218, USA}
\def\Email{\textbf{\color{scured}E-mail}: jgao77@jh.edu}
\def\TEL{\textbf{\color{scured}Phone}: +1 (410)710-0555}
\def\URL{\textbf{\color{scured}Website}: {https://gaojie058.github.io/}}

% ============================================================
% Signature（使用图片签名或文字签名）
% ============================================================
% 如果 signature.pdf 存在，使用图片签名：
% \signature{\vspace{2mm}\includegraphics[scale=0.3]{signature.pdf}\\\vspace{2mm}\name}
% 否则使用文字签名：
\signature{\name}

% ============================================================
% TikZ 坐标系（用于右上角地址块）
% ============================================================
\makeatletter
\def\parsecomma#1,#2\endparsecomma{\def\page@x{#1}\def\page@y{#2}}
\tikzdeclarecoordinatesystem{page}{
    \parsecomma#1\endparsecomma
    \pgfpointanchor{current page}{north east}
    \pgf@xc=\pgf@x%
    \pgf@yc=\pgf@y%
    \pgfpointanchor{current page}{south west}
    \pgf@xb=\pgf@x%
    \pgf@yb=\pgf@y%
    \pgfmathparse{(\pgf@xc-\pgf@xb)/2.*\page@x+(\pgf@xc+\pgf@xb)/2.}
    \expandafter\pgf@x\expandafter=\pgfmathresult pt
    \pgfmathparse{(\pgf@yc-\pgf@yb)/2.*\page@y+(\pgf@yc+\pgf@yb)/2.}
    \expandafter\pgf@y\expandafter=\pgfmathresult pt
}
\makeatother

\address{}

\def\newaddress{
    \Where\\
    \Address\\
    \CityZip\\
    \TEL\\
    \Email\\
    \URL
}

\date{\vspace{10mm} \today}

% ============================================================
% 正文
% ============================================================
\begin{document}

    \begin{letter}{
            \textit{Selection Panel Chair}, \\
            \departmentname{}, \\
            \universityname
        }

        \begin{tikzpicture}[remember picture,overlay,,every node/.style={anchor=center}]
        \node[text width=7cm] at (page cs:0.5,0.73){\small \newaddress};
        \end{tikzpicture}

        \opening{Dear Selection Committee,}

        % --------------------------------------------------------
        % 第一段：开头定调（简洁，直接点明职位和身份）
        % --------------------------------------------------------
        I write to apply for the \href{{{JOB_URL}}}{\textbf{{{POSITION_TITLE}}}} position
        ({{POSITION_REF}}) in \departmentname{} at \textbf{\universityname}.
        I am currently a \href{https://malonecenter.jhu.edu/people-2/members/#postdoctoral-members}{Malone Postdoctoral Fellow}
        at Johns Hopkins University, specializing in \textbf{human-AI collaboration},
        \textbf{responsible and trustworthy AI}, and \textbf{human-centered computing}.
        My research directly aligns with \departmentname's focus in \searchingfocus.

        % --------------------------------------------------------
        % 第二段：研究核心贡献（高影响，具体成果）
        % --------------------------------------------------------
        \textbf{(1) My research addresses fundamental challenges at the intersection of HCI and AI.}
        A central question driving my work is: \textit{How can we design AI systems that strengthen
        human capabilities in analyzing complex unstructured data?} To address this, I have developed
        human-AI partnership systems: \href{https://dl.acm.org/doi/abs/10.1145/3617362}{CoAIcoder}
        (TOCHI 2023) and \href{https://dl.acm.org/doi/abs/10.1145/3613904.3642002}{CollabCoder}
        (CHI 2024), which support efficient, transparent, and trustworthy human-AI analytical workflows.
        Two of these systems---\href{https://gaojie058.github.io/CollabCoder/}{CollabCoder} and
        \href{https://mindcoder.ai/}{MindCoder}---have been publicly deployed as open-source tools,
        demonstrating real-world impact beyond publication. My CHI and TOCHI publications are
        well-cited and have attracted researchers across software engineering, healthcare, and education.

        % --------------------------------------------------------
        % 第三段：与该校 faculty 的具体 overlap（个性化，Step 3 填充）
        % --------------------------------------------------------
        \textbf{(2) Research fit with \universityname{}.}
        {{FACULTY_OVERLAP}}

        % --------------------------------------------------------
        % 第四段：资助潜力（澳洲需要提 ARC）
        % --------------------------------------------------------
        \textbf{(3) Research funding potential.}
        As an early-career researcher with a strong publication record at top-tier venues
        (CHI, TOCHI, UIST), I am well-positioned to pursue competitive funding.
        {{GRANT_MENTION}}
        I also have strong potential to pursue ARC Linkage Projects through collaborations
        with industry partners, building on my prior experience at the Singapore--MIT Alliance
        (SMART), where I collaborated directly with Changi Airport and Synergy Marine Group.

        % --------------------------------------------------------
        % 结尾
        % --------------------------------------------------------
        {{CLOSING_PARA}}

        \closing{Sincerely yours,}
    \end{letter}

\end{document}
